\chapter{Resultados e Discussão}
\label{cap:resultados}

% ─── TODO: preencher conforme os pipelines forem concluídos ──────────────────

\section{Plataforma Implementada}

% Descrever a interface, capturas de tela dos componentes principais.
% Incluir figuras das páginas: dashboard, upload, jobs, PCoA, rede.

\section{Pipeline INOVAHERB}

% Resultados ANCOM-BC2: táxons diferenciais entre tratamentos fatoriais.
% PCoA 1 e PCoA 2 do projeto.
% Rede SpiecEasi: keystone taxa identificados.

\section{Pipeline Pós-Fogo}

% Resultados MaAsLin2: coeficientes de associação por ponto temporal.
% PCoA 3 e PCoA 4 do projeto.
% Predição funcional PICRUSt2.

\section{Pipeline Biorremediação}

% Resultados ANCOM-BC2 (ITS) + FUNGuild: distribuição de guildas.
% Random Forest: variáveis edáficas mais preditivas.
% PCoA 5 e PCoA 6 do projeto.

\section{Painel Unificado de PCoAs}

O painel de seis PCoAs — evento \texttt{CrossProjectFigureReady} da
plataforma — é gerado automaticamente quando os três projetos atingem
o status \texttt{done} em suas análises primárias. A
Figura~\ref{fig:cross_project} apresenta o resultado.

% \begin{figure}[H]
%   \centering
%   \includegraphics[width=\textwidth]{img/cross_project_pcoa.pdf}
%   \caption{Painel unificado com seis PCoAs (duas por projeto):
%            INOVAHERB (a, b), Pós-Fogo (c, d) e Biorremediação (e, f).}
%   \fonte{O autor (2026).}
%   \label{fig:cross_project}
% \end{figure}

\section{Desempenho dos Pipelines}

% Tabela com tempos médios de execução por análise, comparar com estimativas
% do CLAUDE.md (DESeq2: 4–5 min; ANCOM-BC2: 6–8 min; SpiecEasi: 10–20 min...).

\section{Discussão}

% Discutir limitações conhecidas:
% - GSEA hardcoded para org.Hs.eg.db
% - FUNGuild depende de conexão externa
% - Worker single-threaded (jobs longos bloqueiam fila)
% - WebSocket ainda não conectado ao NOTIFY do PG
